\documentclass[12pt,a4paper]{article}
        \makeatletter
        \def\input@path{{../}}
        \makeatother
        \usepackage{almeria-fichas}

        \FichaMeta{Sesion 16}{Relaciones y proporcionalidad}{Bloque 3 · Datos, funciones y graficas}{Reconocer relaciones funcionales y distinguir cuando existe proporcionalidad directa entre dos variables.}{Conceptual}{Comprender, Analizar, Evaluar}

        \begin{document}

        \FichaCabecera

        \begin{materialesbox}
        \begin{itemize}
    \item Ficha impresa
    \item Lapiz
    \item Regla
    \item Calculadora
\end{itemize}
        \end{materialesbox}

        \begin{pasosbox}
        \begin{enumerate}
    \item Lee la pareja de variables.
    \item Comprueba si al multiplicar una variable la otra tambien se multiplica igual.
    \item Si tienes tabla, calcula el cociente $y/x$ cuando se pueda.
    \item No olvides justificar tu conclusion.
\end{enumerate}
        \end{pasosbox}

        \begin{recordatoriobox}
        \begin{itemize}
    \item Hay proporcionalidad directa si $y=kx$.
    \item En una tabla proporcional, el cociente $y/x$ es constante.
    \item No toda relacion funcional es proporcional.
\end{itemize}
        \end{recordatoriobox}

        \begin{apoyobox}
        \begin{itemize}
    \item Haz una columna extra para el cociente $y/x$.
    \item Si el cociente cambia, no hay proporcionalidad directa.
    \item Escribe siempre con que datos has comprobado la relacion.
\end{itemize}
        \end{apoyobox}

        \BloomSection{comprender}

\BloomExercise{comprender}{1}{Que significa proporcional}{Explica con tus palabras que quiere decir que dos variables sean directamente proporcionales.}{6}

\BloomExercise{comprender}{2}{No siempre proporcional}{Explica por que \emph{edad y altura} pueden estar relacionadas pero no ser proporcionales.}{6}

\BloomSection{analizar}

\BloomExercise{analizar}{1}{Dos tablas}{Analiza cual de estas dos tablas es proporcional y cual no. Justifica usando el cociente.}{7}

\BloomExercise{analizar}{2}{Contexto real}{Analiza si \emph{distancia recorrida y tiempo} son proporcionales cuando una persona se para a descansar.}{7}

\BloomSection{evaluar}

\BloomExercise{evaluar}{1}{Modelo correcto}{Evalua si la expresion $y=5x+2$ puede representar una proporcionalidad directa. Justifica.}{6}

\BloomExercise{evaluar}{2}{Decide y argumenta}{Entre varias situaciones de ciudad, evalua cuales incluirias en la guia como ejemplos de proporcionalidad y por que.}{7}

        \begin{evidenciabox}
        \begin{itemize}
    \item Diferenciar relacion funcional y proporcionalidad directa.
    \item Usar tablas para verificar un modelo.
    \item Explicar por que una situacion si o no es proporcional.
\end{itemize}
        \end{evidenciabox}

        \end{document}
